\section{Ensembles and Boosting}
\label{sec:ensembles}

\subsection{Regression trees}

A regression tree partitions the feature space into disjoint regions
$R_1,\dots,R_J$ and predicts the mean target value within each region. A
single split at feature $j$, threshold $s$ is chosen to minimize the total
sum of squared residuals across the two resulting regions:

\begin{equation}
\label{eq:tree-split}
(j^*, s^*) = \arg\min_{j,s}
\left[
  \sum_{i:\, x_i \in R_1(j,s)} \left(y_i - \bar y_{R_1}\right)^2
  +
  \sum_{i:\, x_i \in R_2(j,s)} \left(y_i - \bar y_{R_2}\right)^2
\right]
\end{equation}

Trees require no feature scaling: Equation~\ref{eq:tree-split} depends
only on the ordering of feature values within a column, not their
magnitude, unlike the Ridge objective in Equation~\ref{eq:ridge}.

\subsection{Bagging}

Bagging (bootstrap aggregating) reduces variance by averaging many trees,
each fit to an independent bootstrap resample of the training rows:

\begin{equation}
\label{eq:bagging}
\hat f_{\mathrm{bag}}(x) = \frac{1}{B}\sum_{b=1}^{B} \hat f_b(x)
\end{equation}

where each $\hat f_b$ is a tree fit to bootstrap sample $b$
\cite{breiman1996bagging}. Random forests extend
Equation~\ref{eq:bagging} by additionally restricting each split in
Equation~\ref{eq:tree-split} to a random subset of features, further
decorrelating the individual trees \cite{breiman2001random}.

\subsection{Gradient boosting}

Where bagging averages independent trees fit in parallel, gradient
boosting fits trees sequentially, each one correcting the errors of the
ensemble built so far. The model is an additive expansion

\begin{equation}
\label{eq:gb-additive}
F_M(x) = \sum_{m=1}^{M} \gamma_m h_m(x)
\end{equation}

built greedily: at each stage $m$, tree $h_m$ is fit to the negative
gradient (pseudo-residual) of the loss function $L$ with respect to the
current model's predictions,

\begin{equation}
\label{eq:gb-pseudoresidual}
r_{i,m} = -\left.
\frac{\partial L(y_i, F(x_i))}{\partial F(x_i)}
\right|_{F=F_{m-1}}
\end{equation}

which for squared-error loss reduces to the ordinary residual
$y_i - F_{m-1}(x_i)$. This is gradient descent carried out in the space of
functions rather than parameters \cite{friedman2001greedy}.
\texttt{HistGradientBoostingRegressor} (scikit-learn), used in this
project's Session 4, implements
Equations~\ref{eq:gb-additive}--\ref{eq:gb-pseudoresidual} with
histogram-binned feature values, which reduces the split-search cost of
Equation~\ref{eq:tree-split} from scanning every unique value to scanning
a fixed number of bins, and handles missing feature values natively.

\subsection{XGBoost}
\label{sec:xgboost}

XGBoost augments gradient boosting with an explicit regularization term
and a second-order (Newton) approximation of the loss. Its objective at
boosting round $m$ is

\begin{equation}
\label{eq:xgb-objective}
\mathcal{L}^{(m)} =
\sum_{i=1}^n l\!\left(y_i,\, F_{m-1}(x_i) + h_m(x_i)\right)
+ \Omega(h_m),
\qquad
\Omega(h) = \gamma T + \tfrac{1}{2}\lambda \|w\|^2
\end{equation}

where $T$ is the number of leaves in tree $h_m$ and $w$ its leaf weights.
This project treats XGBoost as a comparison model, not the project's goal
(see \texttt{AGENTS.md}), with a small, defensible hyperparameter grid
selected on the validation years rather than a wide search, implemented
in \texttt{notebooks/05\_xgboost.ipynb} and \texttt{src/train.py}.

\subsubsection{The second-order (Newton) step}

Rather than fitting $h_m$ to the first-order pseudo-residual alone
(Equation~\ref{eq:gb-pseudoresidual}), XGBoost approximates the loss in
Equation~\ref{eq:xgb-objective} by its second-order Taylor expansion
around the current prediction $F_{m-1}(x_i)$. Writing, for each
observation,

\begin{equation}
\label{eq:xgb-grad-hess}
g_i = \left.\frac{\partial\, l(y_i, F)}{\partial F}\right|_{F=F_{m-1}(x_i)},
\qquad
h_i = \left.\frac{\partial^2 l(y_i, F)}{\partial F^2}\right|_{F=F_{m-1}(x_i)},
\end{equation}

(the Hessian $h_i$ carries an observation subscript $i$; the tree $h_m$
carries a round subscript $m$, so the two symbols do not collide), the
objective becomes, up to a constant that does not depend on $h_m$,

\begin{equation}
\label{eq:xgb-taylor}
\tilde{\mathcal{L}}^{(m)} =
\sum_{i=1}^{n}\left[ g_i\, h_m(x_i) + \tfrac{1}{2} h_i\, h_m(x_i)^2 \right]
+ \gamma T + \tfrac{1}{2}\lambda \sum_{j=1}^{T} w_j^2 .
\end{equation}

A tree with fixed structure assigns every observation to exactly one leaf
$j$ with constant value $w_j$. Collecting the per-leaf sums
$G_j = \sum_{i \in I_j} g_i$ and $H_j = \sum_{i \in I_j} h_i$, where
$I_j$ is the set of observations in leaf $j$, Equation~\ref{eq:xgb-taylor}
separates across leaves:

\begin{equation}
\label{eq:xgb-per-leaf}
\tilde{\mathcal{L}}^{(m)} =
\sum_{j=1}^{T}\left[ G_j w_j + \tfrac{1}{2}\left(H_j + \lambda\right) w_j^2 \right]
+ \gamma T .
\end{equation}

Each bracket is a one-dimensional quadratic in $w_j$, minimized in closed
form:

\begin{equation}
\label{eq:xgb-leaf-weight}
w_j^{*} = -\frac{G_j}{H_j + \lambda},
\qquad
\tilde{\mathcal{L}}^{(m)}_{\min} =
-\tfrac{1}{2}\sum_{j=1}^{T}\frac{G_j^2}{H_j + \lambda} + \gamma T .
\end{equation}

\subsubsection{Split gain and pruning}

Because Equation~\ref{eq:xgb-leaf-weight} scores any fixed tree structure,
the value of splitting one leaf into left and right children ($L$, $R$) is
the reduction in the minimized objective, the split gain
\cite{chen2016xgboost}:

\begin{equation}
\label{eq:xgb-gain}
\mathrm{Gain} = \tfrac{1}{2}\left[
\frac{G_L^2}{H_L + \lambda} + \frac{G_R^2}{H_R + \lambda}
- \frac{(G_L + G_R)^2}{H_L + H_R + \lambda}
\right] - \gamma .
\end{equation}

Splits are grown greedily by maximizing Equation~\ref{eq:xgb-gain} over
candidate features and thresholds; any split whose gain is negative after
the $\gamma$ deduction is not made, so $\gamma$ acts as a pruning
threshold on weak splits. (The xgboost implementation reports gain without
the constant factor $\tfrac{1}{2}$, which does not change which split is
chosen; verified against xgboost 3.3.0, whose reported gain is exactly
twice Equation~\ref{eq:xgb-gain} at $\gamma = 0$.)

\subsubsection{Specialization to squared error, and what $\lambda$ does}

Under the squared-error loss $l = \tfrac{1}{2}(y_i - F)^2$ used in this
project, Equation~\ref{eq:xgb-grad-hess} gives
$g_i = F_{m-1}(x_i) - y_i = -r_i$ (the negative residual) and $h_i = 1$,
so for a leaf containing $n_j$ observations Equation~\ref{eq:xgb-leaf-weight}
becomes

\begin{equation}
\label{eq:xgb-shrunk-leaf}
w_j^{*} = \frac{\sum_{i \in I_j} r_i}{n_j + \lambda}
        = \bar r_j \cdot \frac{n_j}{n_j + \lambda},
\end{equation}

where $\bar r_j$ is the mean residual in the leaf. Plain gradient boosting
(Section on gradient boosting above) would output $\bar r_j$ itself;
XGBoost multiplies it by $n_j / (n_j + \lambda)$, shrinking every leaf's
correction toward zero and shrinking thin-evidence leaves (small $n_j$)
hardest. This is the second-order machinery's practical effect on this
project's data: corrections supported by few country-year rows are damped
before they can overfit. Equations~\ref{eq:xgb-leaf-weight},
\ref{eq:xgb-gain}, and~\ref{eq:xgb-shrunk-leaf} were verified numerically
against xgboost 3.3.0 on a worked four-observation example (a single
depth-one tree, $\lambda = 1$): the fitted leaf values equal both
closed-form expressions exactly, and the split chosen is the gain-maximal
one.

\subsubsection{Shrinkage, subsampling, and early stopping}

Each fitted tree enters the ensemble scaled by a learning rate
$\eta \in (0, 1]$:

\begin{equation}
\label{eq:xgb-shrinkage}
F_m(x) = F_{m-1}(x) + \eta\, h_m(x),
\end{equation}

so a smaller $\eta$ requires more boosting rounds $M$ to reach the same
fit; $\eta$ and $M$ trade off directly, and slow learning with more trees
generalizes better at higher compute cost \cite{friedman2001greedy}.
XGBoost additionally fits each tree on a random fraction of rows
(\texttt{subsample}) and columns (\texttt{colsample\_bytree}), which
decorrelates successive trees in the same spirit as bagging
\cite{chen2016xgboost}.

The number of rounds $M$ is not tuned by grid search in this project but
chosen by early stopping: trees are added while the validation-set MAE
improves, and boosting stops when it has not improved for a fixed patience
window, freezing

\begin{equation}
\label{eq:xgb-early-stop}
M^{*} = \arg\min_{m \le M_{\max}} \mathrm{MAE}_{\mathrm{val}}(F_m).
\end{equation}

Early stopping is model selection on the validation years, which this
project's protocol allows; selecting anything by test-set performance is
not. For the single test evaluation, each model is refit on the pooled
training and validation rows (target years through 2018) with the tree
count fixed at $M^{*}$ and early stopping disabled, so the test years
never influence any modeling choice, including the number of trees.

\subsubsection{Why trees require the delta parameterization}

A regression tree predicts a constant within each region of
Equation~\ref{eq:tree-split}, so an ensemble of trees can never predict a
level outside the range of its training targets. A country whose emissions
rise beyond anything seen in training saturates at the training ceiling
and is systematically under-predicted in the level parameterization of
Equation~\ref{eq:level-delta}. Predicting the bounded year-over-year
change $\delta_{c,t+1}$ and reconstructing the level from the known
current value sidesteps this failure mode, which is why the tree-based
models in this project are evaluated in both parameterizations.
